% Cycle10 contribution: exact all-distance response, source-complete proof.
% Authoring agent: /root/cycle10_hamming_response.
% Conventions: P_x projects a full binary input-charge pattern; P_Q pinches
% total input charge. Every statement below is a deep-bulk operator statement.

\section{All-distance coherent response in a weighted input norm}
\label{sec:hammingresponse}

Retain the actual root-retaining device and its ordered $k$-call word.
Write $b=\beta^2$, $a=1-b$, $Q=\sum_iq_i$, $A=Q-k/2$, and
$G_\beta=Q-W_k^*QW_k$. Restrict the initial bath to depths at least $k$.
After total input-charge pinching, this displacement is represented on
the input and its first $k$ bath colors, tensor identity on the seed and
untouched tail. Denote this finite representation by $\mathcal G_k(b)$.
The actual root is retained outside this bulk; it has not been replaced
by a reflecting boundary or a fresh bath.

For $x\in\{0,1\}^k$, let $P_x$ be the complete input charge-pattern
projection and let $\mathcal P(X)=\sum_xP_xXP_x$. All these maps act as
identity on the bath-prefix factor when applied to joint operators.

\begin{theorem}[Uniform all-coherence response]\label{thm:hammingresponse}
Let $\omega$ be any normalized input density with
$\mathcal P(\omega)=\omega$. Define the unital completely positive map
\[
 \Phi_\omega(X)=\operatorname{Tr}_{\rm input}[(\omega\otimes I)X].
\]
For every Hermitian input $Y$ commuting with $Q$, put
$Y_\perp=Y-\mathcal P(Y)$ and
$\|Y_\perp\|_{2,\omega}^2=\operatorname{Tr}(\omega Y_\perp^2)$.
Then, on the full bath-prefix operator space,
\begin{equation}\label{eq:hammingresponse}
 \left\|
 \Phi_\omega[(Y\otimes I)\mathcal G_k(b)]
 -2a\operatorname{Tr}(\omega AY)I
 \right\|
 \le r e^{2r^2}\|Y_\perp\|_{2,\omega},\qquad r=bk.
\end{equation}
In particular the coefficient is at most $2r$ for $0\le r\le1/2$.
For the actual reference
$\sigma=\operatorname{diag}(1/2,1/4,1/4)^{\otimes k}$, every
charge-preserving $Y$ commutes with $\sigma$, so this is the Hermitian
physical signed-input response $D_b(\sigma Y)$.

The exact diagonal identity used here is
\begin{equation}\label{eq:hammingdiagonal}
 \mathcal P(\mathcal G_k(b))=a(2Q-kI)\otimes I.
\end{equation}
The reference may have correlated, nonuniform color fibers. No choice of
bath-prefix state is made in \eqref{eq:hammingresponse}.
\end{theorem}

\subsection{Exact selection at $2m$ flipped endpoints}

Fix input charge patterns $x,y$ of equal weight. Let
$T=\{i:x_i\ne y_i\}$, $|T|=2m$, with $m\ge1$.
There are $m$ zeros and $m$ ones of $x$ on $T$.
At a bulk call, a move flips the binary input charge and a stay retains it.
Thus equality of the two output charge words forces identical choices
at every spectator and complementary choices at each endpoint in $T$.

Let $S\subseteq T^c$ be the common spectator stay set, and
$s\subseteq T$ the ket endpoint stay set. The bra endpoint stay set is
$T\setminus s$. The product of the two amplitudes is
\[
 (ab)^m b^{|S|}a^{k-2m-|S|}\epsilon_s,
 \qquad \epsilon_s=(-1)^{m+|s|}.
\]
Let $R_{s,S}^{yx}$ denote the associated compressed raw cross-history
block, retaining input color wires and all final bath overlaps.
It is a contraction. On this selected charge block the common output
charge has the form
\[
 q_{\rm out}(s,S)=q_{\rm spectator}(x,S)+q_T(x,s),
 \qquad
 q_T(x,s)=m+\sum_{i\in T}(2x_i-1)1_{i\in s}.
\]

Keep spectator calls equal to their fixed $U$ or $I$ choices and endpoint
calls equal to $V_\beta$. This auxiliary ordered word is unitary. Its
$y,x$ block in $W^*W$ vanishes. All its contributing terms have the
same nonzero factor $(ab)^m$ for $0<b<1$, so
\begin{equation}\label{eq:hammingunitarity}
 \sum_{s\subseteq T}\epsilon_sR_{s,S}^{yx}=0.
\end{equation}
The raw blocks do not depend on $b$, so the identity also controls the
endpoint values by continuity. It cancels the entire spectator output
charge, and also permits subtracting the constant $m$ from endpoint
output charge. Since the input $Q$ term is offdiagonal-zero,
\[
 P_y\mathcal G_k(b)P_x
 =(ab)^m\sum_{S\subseteq T^c}
       b^{|S|}a^{k-2m-|S|}B_{yx}^S,
\]
where
\[
 B_{yx}^S=-\sum_{s\subseteq T}
        \epsilon_s(q_T(x,s)-m)R_{s,S}^{yx}.
\]
The spectator weights sum to one. The map from endpoint stay bits to
endpoint output bits is a bijection, hence
\begin{align}
 \|B_{yx}^S\|
 &\le \sum_{j=0}^{2m}\binom{2m}{j}|j-m|
 =m\binom{2m}{m}=:c_m,\notag\\
 \|P_y\mathcal G_k(b)P_x\|&\le c_m(ab)^m.
 \label{eq:hammingblock}
\end{align}
For the displayed binomial identity, pair the terms on the two sides of
$m$ and use
$ (j-m)\binom{2m}{j}
=m[\binom{2m-1}{j-1}-\binom{2m-1}{j}]$
in the upper half; the sum telescopes.
For $m=1$, $c_1=2$, reproducing the completed pair-supported block bound.
For $m\ge2$, all higher-endpoint coherences remain present.

If $x=y$, equality of output charges instead forces the two stay masks
to be identical at every call. For any one mask, the raw word is an
isometry on that input charge pattern, and its output charge is scalar.
The squared amplitudes give mean output charge
$b|x|+a(k-|x|)$. The corresponding displacement is
$|x|-b|x|-a(k-|x|)=a(2|x|-k)$, independently of input colors and the
prefix. This proves \eqref{eq:hammingdiagonal}.

\subsection{Row-square control and operator Cauchy--Schwarz}

Set
\[
 Z=\mathcal G_k(b)-a(2Q-kI)\otimes I.
\]
Thus $Z=Z^*$, $\mathcal P(Z)=0$, and its nonzero input charge blocks
have equal total charge. For fixed $x$ of weight $q$, exactly
$\binom qm\binom{k-q}m$ patterns $y$ have distance $2m$ and the same
weight. Complete-pattern pinching of the square has no cross-column terms:
\[
 P_xZ^2P_x=\sum_y P_xZP_yZP_x.
\]
Each summand is positive and bounded by the square of
\eqref{eq:hammingblock}. Consequently
\begin{align}
 P_xZ^2P_x
 &\le \sum_{m\ge1}\binom qm\binom{k-q}m
                      c_m^2(ab)^{2m}P_x\notag\\
 &\le \sum_{m\ge1}
       \frac{k^{2m}}{4^m(m!)^2}c_m^2 b^{2m}P_x.
 \label{eq:hammingrows}
\end{align}
We used $q(k-q)\le k^2/4$ and $a\le1$; nonexistent binomial terms
are zero. The central binomial bound
$\binom{2m}{m}\le 2^{2m-1}$, valid for every $m\ge1$, now gives
\begin{align}
 \mathcal P(Z^2)
 &\le\sum_{m\ge1}
       \frac{m^24^{m-1}r^{2m}}{(m!)^2}I\notag\\
 &=r^2\sum_{j\ge0}\frac{(4r^2)^j}{(j!)^2}I
 \le r^2e^{4r^2}I.
 \label{eq:hammingrowmoment}
\end{align}
The finite physical sum was only enlarged by the infinite positive series.

Because $\omega=\mathcal P(\omega)$,
$\Phi_\omega(Z^2)=\Phi_\omega(\mathcal P(Z^2))
\le r^2e^{4r^2}I$.
The map $\Phi_\omega$ is unital and completely positive: it can equivalently
be written with $\omega^{1/2}$ on each side of its argument before the
input partial trace. Applying it to the positive $2$-by-$2$ operator matrix
\[
 \begin{pmatrix}
 (Y_\perp^2\otimes I)&(Y_\perp\otimes I)Z\\
 Z(Y_\perp\otimes I)&Z^2
 \end{pmatrix}
 =
 \begin{pmatrix}Y_\perp\otimes I\\Z\end{pmatrix}
 \begin{pmatrix}Y_\perp\otimes I&Z\end{pmatrix}
\]
gives, by its scalar upper-left block,
\[
 \|\Phi_\omega[(Y_\perp\otimes I)Z]\|^2
 \le \operatorname{Tr}(\omega Y_\perp^2)\,
        \|\Phi_\omega(Z^2)\|.
\]
No commutation of $Y_\perp$ and $Z$ is needed. The diagonal part of $Y$
paired with $Z$ has zero input partial trace, and the remaining diagonal
response is exactly $2a\operatorname{Tr}(\omega AY)I$.
This proves \eqref{eq:hammingresponse}. Since $e^{1/2}<2$, its final
$2r$ form holds for $r\le1/2$.

\subsection{A separate flip-set norm interface}

For the actual reference $\sigma$, define
$Y_T=\sum_xP_{x\mathbin{\triangle}T}YP_x$, restricting to equal-charge
pairs. Distinct $T$ are orthogonal in the weighted Hilbert--Schmidt inner
product. If $|T|=2m$, $Y_T$ is supported where $Q_T=m$, whose $\sigma$
mass is $p_m=\binom{2m}{m}/4^m$. Hence
\[
 \|\sigma Y_T\|_1
 \le\sqrt{p_m}\|Y_T\|_{2,\sigma}.
\]
Combining the exact $T$ response with Cauchy--Schwarz over $T$ gives
the occasionally sharper finite-$k$ coefficient
\[
 C_k(b)^2=\sum_{m=1}^{\lfloor k/2\rfloor}
 \binom{k}{2m}p_m c_m^2(ab)^{2m}.
\]
Thus the right side of \eqref{eq:hammingresponse} may also be replaced
by $C_k(b)\|Y_\perp\|_{2,\sigma}$. Its first squared term is
$k(k-1)(ab)^2$. The same central-binomial estimate bounds
$C_k(b)\le re^{2r^2}$. This alternate derivation is not needed for the
general-reference row proof.

\subsection{Necessary scope and receiving use}

The diagonal charge term cannot be omitted. For the actual $\sigma$,
take $Y=A$: $\|A\|_{2,\sigma}=\sqrt{k}/2$, while its exact bulk response
is $D_b(\sigma A)=akI/2$. At fixed small $r=bk$, even the normalized
choice $Y=A/\sqrt{k}$ has bounded weighted norm and response diverging
as $\sqrt{k}$. The new theorem isolates this term instead of hiding it
in an offdiagonal estimate.

The bulk premise is also necessary. At the actual vacuum and $k=1$,
the reference $\sigma=\operatorname{diag}(1/2,1/4,1/4)$ gives depth
drift $a/2$, because an input color pushes and input zero performs the
seed exchange. The claimed bulk diagonal response of this reference
would be zero. Root control must therefore remain in the inherited
operator Lyapunov theorem.

The theorem is an exact signed-input response interface. It supplies no
new bath operation and does not itself prove positive input synthesis,
entropy gain or a communication rate. A receiving positive construction
must independently control its charge excess and its $\|\cdot\|_{2,\sigma}$
offdiagonal remainder, retain the physical root, and invoke the full
arbitrary-reference coding bridge. Spectator weights in the proof arise
only after coherent histories have been paired; they are not a classical
noise replacement for the device.
